
\documentclass{article}
\usepackage{colortbl}
\usepackage{makecell}
\usepackage{multirow}
\usepackage{supertabular}

\begin{document}

\newcounter{utterance}

\centering \large Interaction Transcript for game `eqbench', experiment `eq\_bench', episode 95 with reflexion{-}llama{-}3.1{-}8b{-}instruct.
\vspace{24pt}

{ \footnotesize  \setcounter{utterance}{1}
\setlength{\tabcolsep}{0pt}
\begin{supertabular}{c@{$\;$}|p{.15\linewidth}@{}p{.15\linewidth}p{.15\linewidth}p{.15\linewidth}p{.15\linewidth}p{.15\linewidth}}
    \# & \multicolumn{2}{c}{Player} && \multicolumn{2}{c}{Game Master} \\
    \hline

    \theutterance \stepcounter{utterance}  
    & & & \multicolumn{4}{p{0.6\linewidth}}{
        \cellcolor[rgb]{0.9,0.9,0.9}{
            \makecell[{{p{\linewidth}}}]{
                \texttt{\tiny{[P1$\langle$GM]}}
                \\ 
\texttt{Your task is to predict the likely emotional responses of a character in this dialogue:} \\
\\ 
\texttt{Harvey: You're late, Jess.} \\
\texttt{Jess: I know, Harvey. I'm sorry, I...} \\
\texttt{Harvey: Save it. I've heard it all before. You freelancers, always with the excuses.} \\
\texttt{Jess: It's not an excuse, Harvey. I had a project fall through, and...} \\
\texttt{Harvey: And what? You thought you'd just live here for free?} \\
\texttt{Jess: No, of course not. I just need a little more time.} \\
\texttt{Harvey: Time. Always time with you people. You know, when I was your age, I was running my own auto shop.} \\
\texttt{Jess: I know, Harvey. You've told me.} \\
\texttt{Harvey: And you think you've got it tough? Try losing everything to a crooked business partner.} \\
\texttt{Jess: I didn't say I had it tough, Harvey. I just...} \\
\texttt{Harvey: Just what? Thought you'd take advantage of an old man?} \\
\texttt{Jess: No! That's not it at all. I just...} \\
\texttt{Harvey: You just what? Spit it out, girl!} \\
\texttt{Jess: I just need a little more time, Harvey. That's all.} \\
\texttt{Harvey: And what if I don't have any more time to give, Jess?} \\
\texttt{{[}End dialogue{]}} \\
\\ 
\texttt{At the end of this dialogue, Jess would feel...} \\
\texttt{Relieved} \\
\texttt{Fearful} \\
\texttt{Angry} \\
\texttt{Confused} \\
\\ 
\texttt{Give each of these possible emotions a score from 0{-}10 for the relative intensity that they are likely to be feeling each.} \\
\\ 
\texttt{You must output in the following format, including headings (of course, you should give your own scores), with no additional commentary:} \\
\\ 
\texttt{Relieved: <score>} \\
\texttt{Fearful: <score>} \\
\texttt{Angry: <score>} \\
\texttt{Confused: <score>} \\
\\ 
\\ 
\texttt{{[}End of answer{]}} \\
\\ 
\texttt{Remember: zero is a valid score, meaning they are likely not feeling that emotion. You must score at least one emotion > 0.} \\
\\ 
\texttt{Your answer:} \\
            }
        }
    }
     \\ \\

    \theutterance \stepcounter{utterance}  
    & \multicolumn{4}{p{0.6\linewidth}}{
        \cellcolor[rgb]{0.561,0.737,0.561}{
            \makecell[{{p{\linewidth}}}]{
                \texttt{\tiny{[P1$\rangle$GM]}}
                \texttt{Relieved: 0} \\
\texttt{Fearful: 8} \\
\texttt{Angry: 6} \\
\texttt{Confused: 4} \\
            }
        }
    }
    & & \\ \\

    \theutterance \stepcounter{utterance}  
    & & & \multicolumn{2}{p{0.3\linewidth}}{
        \cellcolor[rgb]{0.9,0.9,0.9}{
            \makecell[{{p{\linewidth}}}]{
                \texttt{\tiny{[GM$|$GM]}}
                \texttt{Relieved: 0} \\
\texttt{Fearful: 8} \\
\texttt{Angry: 6} \\
\texttt{Confused: 4} \\
            }
        }
    }
    & & \\ \\

    \theutterance \stepcounter{utterance}  
    & & & \multicolumn{2}{p{0.3\linewidth}}{
        \cellcolor[rgb]{0.9,0.9,0.9}{
            \makecell[{{p{\linewidth}}}]{
                \texttt{\tiny{[GM$|$GM]}}
                \texttt{Relieved: 0} \\
\texttt{Fearful: 8} \\
\texttt{Angry: 4} \\
\texttt{Confused: 0} \\
            }
        }
    }
    & & \\ \\

    \theutterance \stepcounter{utterance}  
    & & & \multicolumn{2}{p{0.3\linewidth}}{
        \cellcolor[rgb]{0.9,0.9,0.9}{
            \makecell[{{p{\linewidth}}}]{
                \texttt{\tiny{[GM$|$GM]}}
                \texttt{game\_result = LOSE} \\
            }
        }
    }
    & & \\ \\

    \theutterance \stepcounter{utterance}  
    & & \multicolumn{4}{p{0.6\linewidth}}{
        \cellcolor[rgb]{0.95,0.95,0.95}{
            \makecell[{{p{\linewidth}}}]{
                \texttt{\tiny{[P1$\leftrightarrow$P]}}
                \texttt{Post{-}episode Reflexion (not sent to GM):} \\
\texttt{New plan: } \\
\\ 
\texttt{Upon re{-}examining the dialogue, I realized that I focused too much on the surface{-}level conflict between Harvey and Jess, but didn't delve deeper into the emotional dynamics of the conversation. I should have considered the power imbalance between the two characters, with Harvey being the authority figure and Jess being the subordinate. } \\
\\ 
\texttt{In my next attempt, I will pay closer attention to the emotional cues in the dialogue, such as Harvey's aggressive tone and Jess's defensive responses. I will also consider the context of the conversation, including the fact that Jess is a freelancer who is asking for more time, which may be perceived as a threat to Harvey's authority.} \\
\\ 
\texttt{To improve my analysis, I will ask myself questions such as:} \\
\\ 
\texttt{* How does Harvey's past experience of losing everything to a crooked business partner affect his current behavior and emotional state?} \\
\texttt{* How does Jess's tone and language convey her emotional state, and how does it interact with Harvey's emotional state?} \\
\texttt{* What are the underlying power dynamics at play in this conversation, and how do they influence the emotional responses of the characters?} \\
\\ 
\texttt{By considering these factors, I hope to provide a more nuanced and accurate analysis of the emotional responses of the characters in the dialogue.} \\
            }
        }
    }
    & \\ \\

\end{supertabular}
}

\end{document}
